\section{Formalisme Teori Permainan}

\begin{frame}
\centering
\highlight{\Huge Bagian 4: Formalisme Teori Permainan}
\end{frame}

\subsection{Utilitas Individu}
\begin{frame}{4.1 Utilitas Individu \& Payoff Pemain}
Dalam model ini, setiap aset $i$ dianggap sebagai pemain yang mencoba memaksimalkan utilitasnya sendiri ($u_i$) berdasarkan keputusan alokasi $x_i \in \{0, 1\}$.

\textbf{Fungsi Utilitas Strategis:}
\begin{equation}
u_i(x_i, x_{-i}) = x_i \underbrace{\left( \tilde{\mu}_i - \frac{\gamma}{2}\tilde{\sigma}_{ii} - \gamma \sum_{j \neq i} \tilde{\sigma}_{ij} x_j \right)}_{\text{Marginal Utility}}
\end{equation}

\textbf{Komponen Payoff:}
\begin{itemize}
    \item $\tilde{\mu}_i$: Imbal hasil strategis (dorongan positif).
    \item $\frac{\gamma}{2}\tilde{\sigma}_{ii}$: Penalti risiko individu (disinsentif variansi).
    \item $\gamma \sum_{j \neq i} \tilde{\sigma}_{ij} x_j$: Penalti interaksi (disinsentif korelasi/redundansi dengan aset lain yang dipilih).
\end{itemize}
\end{frame}

\subsection{Exact Potential Game}
\begin{frame}{4.2 Definisi \textit{Exact Potential Game}}
Sistem ini diklasifikasikan sebagai \textbf{\textit{Exact Potential Game}} karena setiap perubahan utilitas marginal pemain identik dengan perubahan pada sebuah fungsi potensial global ($\Phi$).

\textbf{Syarat Formal:}
\begin{equation}
u_i(1, x_{-i}) - u_i(0, x_{-i}) = \Phi(1, x_{-i}) - \Phi(0, x_{-i})
\end{equation}

\textbf{Signifikan Fisika-Finansial:}
\begin{itemize}
    \item Egoisme pemain (maksimalisasi $u_i$) secara otomatis menuntun sistem menuju optimasi global ($\Phi$).
    \item Titik \textit{Nash Equilibrium} murni dari permainan ini berhimpit dengan titik optimal global (minimum/maksimum) dari fungsi potensial.
\end{itemize}
\end{frame}

\subsection{Konstruksi Fungsi Potensial}
\begin{frame}{4.3 Konstruksi Fungsi Potensial Global ($\Phi$)}
Fungsi potensial $\Phi$ merangkum seluruh dinamika interaksi antar pemain ke dalam satu skalar tunggal. Ini adalah representasi matematika dari "energi" sistem portofolio.

\textbf{Formulasi Potensial:}
\begin{equation}
\Phi(x) = \sum_{i=1}^N x_i \tilde{\mu}_i - \frac{\gamma}{2} \sum_{i=1}^N \sum_{j=1}^N \tilde{\sigma}_{ij} x_i x_j
\end{equation}

\textbf{Analogi Teoretis:}
\begin{itemize}
    \item Suku pertama ($\sum x_i \tilde{\mu}_i$) setara dengan pengaruh medan magnet luar pada spin.
    \item Suku kedua ($\sum \sum \tilde{\sigma}_{ij} x_i x_j$) setara dengan interaksi antar-spin (kopling) dalam model Ising.
\end{itemize}
\end{frame}

\subsection{Interaksi Strategis}
\begin{frame}{4.4 Interaksi Strategis \& Penalti Marginal}
Inti dari \textit{Potential Game} adalah adanya hukuman bersama jika terjadi redundansi informasi dalam pemilihan aset.

\textbf{Mekanisme Penalti:}
\begin{itemize}
    \item Jika aset $i$ dan $j$ memiliki korelasi informasi tinggi ($\text{NMI}_{ij} \gg 0$), maka $\tilde{\sigma}_{ij}$ akan bernilai besar.
    \item Saat keduanya dipilih ($x_i=1, x_j=1$), utilitas keduanya berkurang secara signifikan sebesar $\gamma \tilde{\sigma}_{ij}$.
\end{itemize}

\textbf{Kesimpulan:}
Sistem secara cerdas akan menghindari aset yang "serupa" (redundant) dan lebih memilih diversifikasi yang memiliki komplementaritas informasi tinggi untuk meminimalkan penalti pada $\Phi(x)$.
\end{frame}
